\subsection{Graded resolutions}

In this section, one supposes that the ring A is equipped with a grading \((\mathbf{A}_n)_{n\in \mathbf{Z}}\) , such that \(\mathbf{A}_n = 0\) for \(n < 0\) . One says that a graded A-module M is bounded below if \(\mathbf{M}_n = 0\) for n sufficiently small; every finitely generated graded A-module is bounded below.

PROPOSITION 11. --- If M is a graded A-module bounded below (resp. if M is a finitely generated graded A-module and if A is left Noetherian), there exists an exact sequence of graded A-modules, unbounded to the left,

\[
\dots \longrightarrow \mathrm{L} _ {n} \xrightarrow {d _ {n}} \mathrm{L} _ {n - 1} \longrightarrow \dots \longrightarrow \mathrm{L} _ {1} \xrightarrow {d _ {1}} \mathrm{L} _ {0} \xrightarrow {d _ {0}} \mathrm{M} \longrightarrow 0
\]

where the \(L_{i}\) are graded free and bounded below (resp. graded free and finitely generated), and where the \(d_{i}\) are graded homomorphisms of degree 0.

If N is a graded A-module bounded below (resp. and finitely generated over Noetherian A) there exists a graded free A-module bounded below (resp. and finitely generated) L and a surjective graded homomorphism L → N (II, p.~167, remark 3).

That being so, suppose given an exact sequence of graded A-modules and graded homomorphisms of degree 0

\[
\mathrm{L} _ {n} \xrightarrow {d _ {n}} \mathrm{L} _ {n - 1} \longrightarrow \dots \longrightarrow \mathrm{L} _ {0} \xrightarrow {d _ {0}} \mathrm{M} \longrightarrow 0,
\]

where the \(L_{i}, i = 0, \ldots, n\) , are graded free and bounded below (resp. graded free and finitely generated). Then N = Ker \(d_{n}\) is bounded below (resp. finitely generated); there therefore exists a graded free A-module bounded below (resp. graded free and finitely generated) \(L_{n+1}\) and a graded homomorphism \(d_{n+1}: L_{n+1} \to L_{n}\) of degree 0, such that Im \(d_{n+1} = N\) ; the sequence

\[
\mathrm{L} _ {n + 1} \xrightarrow {d _ {n + 1}} \mathrm{L} _ {n} \xrightarrow {d _ {n}} \mathrm{L} _ {n - 1} \longrightarrow \dots \longrightarrow \mathrm{L} _ {0} \xrightarrow {d _ {0}} \mathrm{M} \longrightarrow 0
\]

is then exact. The proposition then follows from the preceding by induction on n.

